\section{Discussion}
\label{sec:discussion}

\subsection{Why focal alone wins: the class-imbalance hypothesis}
\label{sec:discussion:imbalance}

Our headline result — focal loss in plain source-only training
outperforms all adversarial methods on WASSA-21 LODO — is consistent
with a simple but unstated hypothesis: \emph{the dominant cross-dataset
shift in emotion is class-distribution shift, not feature-distribution
shift.}

Cross-dataset emotion benchmarks combine corpora with materially
different class frequencies and, in the case of ISEAR, an outright
missing class. Table~\ref{tab:datastats} quantifies the asymmetry:
GoEmotions is 38:1 imbalanced toward \emph{joy} (19{,}763 examples),
ISEAR is uniformly balanced across the five Ekman classes it
contains, and WASSA-21 has \emph{joy} as its \emph{rarest} class
(99 examples, 5.2\% of the corpus). Under LODO with WASSA-21 as the
target, the source training pool inverts this entirely: 51\% of
training examples are labeled joy, while the held-out target has joy
as its rarest class (9 of 164 examples). This is a textbook example
of class-distribution shift across domains; it is precisely the
condition that focal loss \citep{lin2017focal} was designed to
address.

Adversarial methods (DANN, CDAN) are designed for the opposite
problem: stabilizing feature representations across domains that
share the label distribution but differ in input characteristics.
When the actual problem is that the source domains overrepresent
certain classes and the held-out target distributes labels
differently, adversarial alignment cannot help: it would require the
discriminator to enforce that the source-domain representations of
"joy" look like the target-domain representations of "joy", but the
source-domain representations of joy are themselves biased toward
the source's overrepresented uses of that label (e.g.\ GoEmotions
Reddit positivity vs WASSA empathy essays).

Focal loss directly attacks the class-imbalance problem during
training: by down-weighting well-classified examples and concentrating
gradient on rare or hard examples, it produces a classifier less
biased toward the majority classes of the source training set.
Section~\ref{sec:results:perclass} provides the per-class evidence:
focal recovers the two classes (disgust, fear) where CE source-only
collapses to near-zero recall, at a small expense on a class
(surprise) where CE was already doing well.

\subsection{Why adversarial + focal does not amplify}
\label{sec:discussion:combinations}

On WASSA-21 LODO, the three focal-based methods rank as: CE + Focal
($0.516$) $>$ DANN + Focal ($0.501$) $>$ CDAN + Focal ($0.483$).
\emph{Adding} an adversarial layer to focal training reduces the
gain.

Two interpretations are consistent with this pattern. First, the
gradient capacity is finite: the adversarial loss term competes with
the focal regularizer for gradient direction. When the dominant
shift is class-distribution, the adversarial gradient is at best
neutral and at worst aligned with the source-overrepresented
classes — pulling the model back toward the very imbalance that
focal was correcting. Second, adversarial alignment on
\emph{source-source} pairs (as in our LODO regime, where the
discriminator sees only two source domains during training) does not
extrapolate to the held-out target. CDAN, by conditioning the
discriminator on task-head predictions, tightens this source-source
alignment further; this is consistent with CDAN's stronger negative
effect than DANN's in the combinations.

\subsection{Two failure modes for adversarial methods}
\label{sec:discussion:failure}

We characterize two distinct failure modes in our experiments.

\paragraph{Mode 1 — DANN instability at moderate $\lambda$.}
At $\lambda_{\max}=1.0$ (the default in \citet{ganin2016dann}),
domain-loss diverges around epoch 3 of training on all three LODO
targets. On GoEmotions LODO seed 42, the domain loss rises from 1.7
to 10+ and never recovers, while task loss
stays elevated. We mitigated this by halving $\lambda_{\max}$ to
$0.5$ in the main results, but even at the stabilized value DANN
does not exceed CE source-only on any LODO target. We provide the
lambda trajectory in
Figure~\ref{fig:lambda} and the underlying logs in
Appendix~\ref{app:lambda}.

\paragraph{Mode 2 — CDAN source-manifold tightening.}
CDAN training is numerically stable across all LODO targets and seeds,
unlike DANN. But on GoEmotions LODO — the target with the largest
register distance from the source training data — CDAN produces test
macro-F1 essentially indistinguishable from CE
($0.294 \pm 0.015$ vs $0.300 \pm 0.007$). Source validation reaches
$0.77$ at the best epoch, but the gap to test ($0.47$ F1 points) is
the largest of any method-target combination. CDAN's conditional
alignment tightens the source-source feature manifold; when the
held-out target lies far outside that manifold, the tightened
representation does not generalize.

\begin{figure}[t]
\centering
\includegraphics[width=0.95\linewidth]{../outputs/figures/dann_lambda_dynamics.pdf}
\caption{Per-epoch domain loss (right panel) and the sigmoid
$\lambda$ schedule (left panel) for DANN training on three LODO
targets at $\lambda_{\max}=1.0$. GoEmotions exhibits unbounded
domain-loss divergence beginning at epoch 3; ISEAR and WASSA-21
remain stable. Reducing $\lambda_{\max}$ to $0.5$ stabilizes all
three targets but does not produce a positive test F1 gain.}
\label{fig:lambda}
\end{figure}

\subsection{Two LODO targets are method-invariant}
\label{sec:discussion:invariant}

On GoEmotions LODO, all six methods produce test macro-F1 in the
narrow band $[0.290, 0.308]$; on ISEAR LODO, in
$[0.481, 0.517]$. No pair of methods reaches $p < 0.1$ on either
target. We interpret this method-invariance structurally rather
than statistically.

GoEmotions consists of short Reddit comments with sarcasm, slang, and
context-dependent emotion expression — a register essentially absent
from the LODO source training pool (ISEAR self-reports plus WASSA
empathy essays, both long-form first-person reflective text). When
the held-out target lies entirely outside the source manifold, no
feature-space adaptation method we tested can recover the target
distribution, because the source training data provides no signal
about target-specific features.

ISEAR's structural barrier is different and more severe: the dataset
contains zero \emph{surprise} examples. A model trained on
GoEmotions plus WASSA (both of which contain surprise) cannot learn
the target distribution for that class, because it has no
target-distribution examples to align to. Macro-F1 on ISEAR is
therefore computed over five Ekman classes (restrict-to-present), and
the method's score is bounded by how well it learned the other five
on training data that overrepresents goemotions-style joy and ISEAR-style
sadness.

These observations are not weaknesses of our experimental design;
they are structural facts about the cross-dataset emotion benchmark
that the field has tolerated implicitly. Our explicit reporting of
the method-invariance pattern is itself a contribution: it sets a
realistic floor on what feature-space cross-dataset methods can be
expected to deliver on heterogeneous emotion corpora.

\subsection{Scale: why the DeBERTa-large result is variance, not signal}
\label{sec:discussion:scale}

The five-seed DeBERTa-large run on the same cell (CE + Focal $\times$
WASSA-21 LODO) produces a mean of $0.492$ with standard deviation
$0.029$, compared to base mean $0.516$ with standard deviation
$0.005$ at three seeds. Mean trends downward; variance is six times
larger.

We attribute the variance increase to the interaction between
WASSA-21's small target test set ($n = 164$ across six classes,
median 17 per class) and the larger model's higher capacity. Macro-F1
on 16-class buckets is sensitive to a small number of misclassifications;
a high-capacity model can learn fine-grained
source-domain features that, at some random initializations, generalize
well to the target and at others do not. Source-validation macro-F1
remains tightly clustered ($[0.677, 0.683]$, std $0.002$), evidence
that the variance is downstream of training stability, not training
itself.

This result does not undermine the base-model contribution but does
clarify the regime: focal loss helps cross-dataset emotion under
moderate capacity and ample relative test size; with larger capacity
and small test sets the gain becomes less measurable. We report this
honestly rather than scaling the seed count until significance is
recovered.
